\documentclass[12pt]{article}
\usepackage[utf8]{inputenc}
\usepackage[spanish]{babel}
\usepackage{amsmath, amssymb, amsthm}
\usepackage{tcolorbox}
\usepackage{geometry}
\geometry{margin=2.5cm}

\title{Matemática Discreta \\[0.5em]
\large El sistema de los enteros y Anillos y aritmética modular}
\author{Matemática Discreta - UADER 2025}
\date{}

\begin{document}
\maketitle
\tableofcontents
\bigskip
\newpage


\section{Unidad 5: El sistema de los enteros}

\subsection*{Algoritmo de la división: números primos}

\begin{tcolorbox}[title=Algoritmo de la división]
    \begin{itemize}
        \item Sean \( a, b \in \mathbb{Z} \Rightarrow \exists ! q, r \in \mathbb{Z} : \)
            \[ a = qb + r, \quad 0 \leq r < b \] 
            Donde: \( a \) es el \textbf{dividendo}, \( b \) es el \textbf{divisor}, \( q \) es el \textbf{cociente} y \( r \) es el \textbf{resto}. Si \( r = 0 \), entonces \( a \mid b \).
        \item Si \( a, b \in \mathbb{Z} \) y \( b \neq 0 \), entonces \( b \mid a \) si y solo si \( \exists n \in \mathbb{Z} \) tal que \( a = bn \).
        %Agregar nota: se dice que "b divida a a", "b es un divisor de a" o "a es múltiplo de b"
    \end{itemize}
\end{tcolorbox}

\begin{tcolorbox}[title=Propiedades de la relacion de divisibilidad]
\begin{enumerate}
    \item \( 1 \mid a \) (1 es el \textit{divisor universal})
    \item \( a \mid 0 \)
    \item \( (a \mid b \land b \mid a) \Rightarrow a = \pm b \)
    \item \( (a \mid b \land b \mid c) \Rightarrow a \mid c \)
    \item \( a \mid b \Rightarrow a \mid bx \)
    \item Si \( x = y + z \land a \text{ divide a dos de } x, y, z \Rightarrow a \text{ divide al restante} \) 
    \item \( (a \mid b \land a \mid c) \Rightarrow a \mid (bx + cy) \)
    %Agregar nota de que se combinan las propiedades 5 y 6
\end{enumerate}
\end{tcolorbox}

%Agregar nota de definicionnde numero primo y compuesto

\begin{tcolorbox}[title=Número primos y compuestos]
    Si \( n \) es compuesto, \( \Rightarrow \exists ! \text{ un primo } p : p \mid n \land p \leq \sqrt{n} \)
\end{tcolorbox}
% Agregar el contrarrecíproco

\subsection*{Máximo común divisor}

\begin{tcolorbox}[title=Algoritmo de la división]

\end{tcolorbox}

%------------------------
\newpage

\section{Unidad 6: Anillos y aritmética modular}

\subsection*{Estructura de un anillo}

\begin{tcolorbox}[title=Definición de anillo]
Sea \( R \) un conjunto no vacío con dos operaciones binarias cerradas, denotadas por \( + \) y \( \cdot \), que pueden ser diferentes de las operaciones usuales. El conjunto \( (R, +, \cdot) \) es un \textbf{anillo} si para todos \( a, b, c \in R \) se cumplen las siguientes condiciones:
\begin{enumerate}
    \item \( R \neq \emptyset \) (\textbf{conjunto no vacío})
    \item \( a + b \in R \) (\textbf{cerradura} de la adición)
    \item \( a \cdot b \in R \) (\textbf{cerradura} de la multiplicación)
    \item \( a + b = b + a \) (\textbf{conmutativa} de la adición)
    \item \( a + (b + c) = (a + b) + c \) (\textbf{asociativa} de la adición)
    \item \( a \cdot (b \cdot c) = (a \cdot b) \cdot c \) (\textbf{asociativa} de la multiplicación)
    \item \( \exists z \in R : a + z = z + a = a \) (\textbf{identidad aditiva o neutro}, o de la adición)
    \item \( \forall a \in R, \exists (-a) \in R : a + (-a) = (-a) + a = z \) (\textbf{inverso aditivo}, o de la adición)
    \item \( w \cdot (x + y) = w \cdot x + w \cdot y \) (\textbf{distributiva} de la multiplicación respecto a suma a la izquierda)
    \item \( (x + y) \cdot w = x \cdot w + y \cdot w \) (\textbf{distributiva} de la multiplicación respecto a suma a la derecha)
\end{enumerate}
\end{tcolorbox}

\begin{tcolorbox}[colback=yellow!15, title=Aclaración]
\textbf Se hace la aclaración ``a la izquierda'' y ``a la derecha'' porque en las matrices la multiplicación no es conmutativa.
\end{tcolorbox}


\subsection*{Propiedades adicionales}

\begin{tcolorbox}[title=Propiedades importantes]
En cualquier anillo \( (R, +, \cdot) \):
\begin{itemize}
    \item Si \( ab = ba, \forall a, b \in R \), el anillo es \textbf{conmutativo}.
    \item El anillo tiene \textbf{divisores propios de cero} si \( \exists a, b \in R, a \neq z, b \neq z : ab = z \).
    \item Si \( \exists u \in R : u \cdot a = a \cdot u = a, \forall a \in R \), \(\boldsymbol{u}\) es el elemento \textbf{unidad} o identidad para la multiplicación.
    \item Si el anillo tiene elemento unidad, un elemento \( a \in R \) es una \textbf{unidad} si \( \exists b \in R : ab = ba = u \). En este caso, \( b \) es el \textbf{inverso multiplicativo} de \( a \).
    \item Un \textbf{dominio de integridad} es un anillo conmutativo, con elemento unidad, y \textbf{sin divisores propios de cero}.
    \item Un \textbf{cuerpo} es un anillo conmutativo y con elemento unidad en el que todo elemento diferente de cero tiene \textbf{inverso multiplicativo}. Si un anillo es un cuerpo, entonces es un dominio de integridad (por ende no posee divisores propios de cero).
\end{itemize}
\end{tcolorbox}

\begin{tcolorbox}[colback=yellow!15, title=Nota]
\begin{itemize}
    \item El \textbf{neutro}, el \textbf{inverso aditivo} de cada elemento y el elemento \(\boldsymbol{u}\), si existen, son únicos. 
\end{itemize}
\end{tcolorbox}

\begin{tcolorbox}[title=Leyes de cancelación para la suma y absorción]
\begin{itemize}
    \item  \( a + b = b + c = b + a = c + a \Rightarrow b = c \) (\textbf{cancelación para la suma})
    \item \( a \in R : za = az = z \) (\textbf{absorción})
\end{itemize}
\end{tcolorbox}

\newpage

\subsection*{Subanillos e ideales}

\begin{tcolorbox}[title=Subanillo]
Sea \( (R, +, \cdot) \) un anillo y \( \emptyset \neq S \subseteq R \). Decimos que S es un \textbf{subanillo} de R si cumple:
\begin{itemize}
    \item En anillos \textbf{infinitos}: \( \forall a,b \in S, \boldsymbol{a - b} \in S \land ab \in S \)
    \item En anillos \textbf{finitos}: \( \forall a,b \in S, \boldsymbol{a + b} \in S \land ab \in S \)
\end{itemize}
\end{tcolorbox}

\begin{tcolorbox}[title=Ideales]
Un subconjunto no vacío \( I \subseteq R \) es un \textbf{ideal} si:
\begin{itemize}
    \item \( \forall a, b \in I, a - b \in I \).
    \item \( \forall r \in R, ra \land ar \in I \).
\end{itemize}
\end{tcolorbox}

\begin{tcolorbox}[colback=yellow!15, title=Importante]
Todo ideal es un subanillo, pero \textbf{no todos los subanillos son ideales}. La diferencia radica en la cerradura multiplicativa del lado del elemento externo \( r \).
\end{tcolorbox}

\newpage

\subsection*{Los enteros módulo \( n \)}

\begin{tcolorbox}[title=Congruencia modular]
Sea \( n \in \mathbb{Z}^+, n > 1 \). Para \( a, b \in \mathbb{Z} \), decimos que
\[
a \equiv b \pmod{n}
\]
si \( n \mid (a - b) \), es decir, si \( a - b = kn \) para algún \( k \in \mathbb{Z} \).

\bigskip

\textbf{Nota:} esta relación de equivalencia divide a \( \mathbb{Z} \) en \( n \) clases de equivalencia.
\end{tcolorbox}

\begin{tcolorbox}[title=Relación de equivalencia]
\begin{itemize}
    \item \textbf{Reflexiva:} \( a \equiv a \pmod{n} \)
    \item \textbf{Simétrica:} Si \( a \equiv b \pmod{n} \Rightarrow b \equiv a \pmod{n} \)
    \item \textbf{Transitiva:} Si \( a \equiv b \pmod{n} \land b \equiv c \pmod{n} \Rightarrow a \equiv c \pmod{n} \)
\end{itemize}
\end{tcolorbox}

\begin{tcolorbox}[title=Anillo de enteros \( \pmod{n} \)]
Para \( n \in \mathbb{Z}^+, n > 1 \), \( \mathbb{Z}_n \) es un anillo finito y conmutativo, con elemento unidad igual a [1] en las operaciones binarias:

\bigskip

En términos de clases: \( \forall [a], [b] \in \mathbb{Z} \)
\begin{align*}
    [a]_n + [b]_n &= [a + b]_n \\
    [a]_n \cdot [b]_n &= [a b]_n
\end{align*}
y en términos de módulo: si \( a \equiv b \pmod{n} \) y \( c \equiv d \pmod{n} \), entonces:
\begin{align*}
    a + c &\equiv b + d \pmod{n} \\
    ac &\equiv bd \pmod{n}   
\end{align*}
\end{tcolorbox}

\begin{tcolorbox}[colback=yellow!15, title=Nota]

Si \( a \equiv b \pmod{n} \), entonces sus clases de equivalencia son iguales: \( [a] = [b] \) en \( \mathbb{Z}_n \).
\end{tcolorbox}

\begin{tcolorbox}[title=Ejemplo de clases de equivalencia]
En \( \mathbb{Z}_6 \), la clase del 1, denotada \( [1] \), es el conjunto de todos los enteros que son congruentes con 1 módulo 6.
\[ [1] = \{ \dots, -11, -5, 1, 7, 13, \dots \} \]
Como 7 y 13 están en este conjunto, podemos decir que \( 7 \equiv 1 \pmod{6} \) y \( 13 \equiv 1 \pmod{6} \).

\bigskip

Esto significa que las clases \( [1] \), \( [7] \) y \( [13] \) son exactamente el mismo conjunto, por lo que podemos escribir:
\[ [1] = [7] = [13] \]
\end{tcolorbox}

\bigskip

\begin{tcolorbox}[title=Propiedades del anillo \( \mathbb{Z}_n \)]
\begin{itemize}
    \item \( \mathbb{Z}_n \) es un \textbf{cuerpo} \( \boldsymbol{\Leftrightarrow n} \)\textbf{es primo}.
    \item Si \( \exists [a] \in \mathbb{Z}_n : \text{mcd}(a, n) = 1 \rightarrow [a] \) es una \textbf{unidad} (por ende, tiene inverso multiplicativo).
    \item Hay \( \varphi(n) \) \textbf{unidades} y \( n - 1 - \varphi(n) \) \textbf{divisores propios de cero}.
\end{itemize}
\end{tcolorbox}

\begin{tcolorbox}[title=Ecuación lineal de congruencia]
La ecuación lineal de congruencia \( ax \equiv b \pmod{n} \) tiene solución \( \Leftrightarrow \text{mcd}(a, n) \mid b \).
Si \( d = \text{mcd}(a, n) \rightarrow\) hay \( d \) soluciones \( \neq \text{en} \pmod{n} \)
\end{tcolorbox}

\begin{tcolorbox}[title=Notas]
\begin{itemize}
    \item En \( \mathbb{Z}_n, [1] \equiv 1 \) y \( [n - 1] \equiv -1 \equiv (n - 1) \), y ambos son \textbf{siempre unidades}. 
    % Agregar nota del -1
    \item Sea \( n = p_1^{e_1} p_2^{e_2} \cdots p_k^{e_k} \) la factorización en primos de \( n \), ningún número que contenga alguno de los \( p_i \) como factor es una unidad en \( \mathbb{Z}_n \).
\end{itemize}
\end{tcolorbox}

\newpage

\subsection*{Homomorfismo e isomorfismo}

\begin{tcolorbox}[title=Definición]
Sean \( (R, +, \cdot) \) y \( (S, \odot, \otimes) \) anillos. Una función \( f: R \to S \) es un \textbf{homomorfismo de anillos} si:
\begin{align*}
    f(a + b) &= f(a) \oplus f(b) \\
    f(a b) &= f(a) \odot f(b)
\end{align*}
Si además \( f \) es biyectiva, entonces es un \textbf{isomorfismo}.
Para ser biyectiva debe ser \textit{inyectiva y sobreyectiva}.
\end{tcolorbox}

\begin{tcolorbox}[colback=yellow!15, title=Nota]
    Al relacionar un conjunto infinito con uno finito, no hay isomorfismo, pues no es \textbf{inyectiva}.
\end{tcolorbox}

\newpage

\subsection*{El Teorema Chino del Resto}

\begin{tcolorbox}[title=Isomorfismo]
Sea \( n = p_1^{e_1} \cdots p_k^{e_k} \) la factorización en primos de \( n \) y \( m_1 = p_1^{e_1}, \dots, m_k = p_k^{e_k} \). Entonces:
\[
\mathbb{Z}_n \cong \mathbb{Z}_{m_1} \times \mathbb{Z}_{m_2} \times \cdots \times \mathbb{Z}_{m_k}
\]
\end{tcolorbox}

\begin{tcolorbox}[title=Teorema Chino del Resto]
Sean \(m_1, \dots, m_r \in \mathbb{N}\) \textbf{mutuamente primos} y sea \(M\) su producto.
\( \forall a_1, \dots, a_r \in \mathbb{Z} \), \( \exists ! x \in \mathbb{Z}, \pmod{M} : x \equiv a_i \pmod{m_i}, \forall i = 1, \dots, r \).
\end{tcolorbox}

\begin{tcolorbox}[colback=yellow!15, title=Nota]
\begin{align*}
    M = \prod_{i=1}^r m_i && M_i = \frac{M}{m_i} && v_i \text{ es el inverso de } M_i \pmod{m_i}
\end{align*}
\end{tcolorbox}

\begin{tcolorbox}[title=Teorema de Fermat]    
Si \( p \) es un número primo y \( a \in \mathbb{Z} : p \nmid a \), entonces:
\[
a^{p-1} \equiv 1 \pmod{p}
\]
\end{tcolorbox}

\begin{tcolorbox}[title= Indicador de Euler]    
Si \( m \in \mathbb{N} \), el indicador de Euler \( \varphi(m) \) :
\[ \varphi(m) = |\{k \in \mathbb{N} \mid 1 \le k \le m \land \text{mcd}(k, m) = 1\}| \]
\[
    \varphi(m) = m \cdot \frac{p_1 - 1}{p_1} \cdot \frac{p_2 - 1}{p_2} \cdots \frac{p_r - 1}{p_r}
\]
\end{tcolorbox}

\begin{tcolorbox}[title= Teorema de Fermat-Euler]    
Sea \( m \in \mathbb{N} \) y \( a \in \mathbb{Z} : \text{mcd}(a, m) = 1\), entonces:
\[
a^{\varphi(m)} \equiv 1 \pmod{m}
\]
\end{tcolorbox}

\end{document}
